\section{Research Gaps}
\textbf{Limited integration of modern AI models with Indian datasets}  

While global research uses advanced models like transformers and graph neural networks, most Indian datasets remain underutilized due to size, format, and availability challenges. There is a lack of well-tested frameworks that combine these models effectively for crops grown in Indian conditions.
\\\\
\textbf{Insufficient edge-ready and offline solutions}  

A large portion of Indian farmers lack consistent internet access, which makes many cloud-based systems impractical. Current solutions rarely provide edge or SMS-based communication for farmers in low-connectivity regions. Developing lightweight models that can run on basic hardware is still a major research need.
\\\\
\textbf{Lack of localized precision systems in Indian agriculture}  

Most studies and government platforms focus on large-scale or regional data aggregation, but few target the local village or farm level. Farmers still rely on approximate yield estimates instead of real, data-backed predictions. This gap highlights the need for systems that use real-time, localized environmental and soil data to improve decision-making.
\\\\
\textbf{Data privacy and decentralized learning challenges}  

Federated learning and privacy-preserving AI are still new in Indian agriculture. There are few real-world deployments proving their effectiveness in field conditions. More research is required to build decentralized systems that protect farmer data while maintaining high accuracy.
\\\\
\textbf{Incomplete handling of missing and inconsistent sensor data}  

Field sensors and satellite inputs often have missing or faulty readings due to hardware or weather issues. Many existing systems simply drop incomplete data, leading to poor model reliability. Better imputation and data-repair techniques need to be explored to make yield prediction models robust and trustworthy.
